% =====================================================================
%  CARNET D'ENTRAÎNEMENT — Fiche 12 (Chimie organique)
%  THÈME 4 — Masse molaire d'une molécule organique
%  NIVEAU FACILE — CORRIGÉ
% =====================================================================
\documentclass[11pt,a4paper]{article}
\usepackage[utf8]{inputenc}
\usepackage[T1]{fontenc}
\usepackage{lmodern}
\usepackage[french]{babel}
\usepackage{amsmath,amssymb,mathtools}
\providecommand{\degree}{^\circ}
\usepackage{icomma}
\usepackage{siunitx}
\sisetup{output-decimal-marker={,},per-mode=symbol}
\usepackage[version=4]{mhchem}
\usepackage{chemfig}
\setchemfig{atom sep=2em,bond length=0.9em,cram width=2.5pt}
\usepackage{xcolor}
\usepackage{array}
\usepackage{booktabs}
\usepackage{tabularx}
\usepackage[most]{tcolorbox}
\usepackage{enumitem}
\usepackage{tikz}
\usetikzlibrary{arrows.meta,positioning,calc}
\usepackage{fancyhdr}
\usepackage{titlesec}
\usepackage[hidelinks]{hyperref}
\usepackage[a4paper,margin=2.1cm,headheight=26pt,headsep=9pt]{geometry}

\definecolor{primary}{RGB}{150,98,162}
\definecolor{primarydark}{RGB}{92,52,110}
\definecolor{excol}{RGB}{0,135,90}
\definecolor{attncol}{RGB}{200,40,55}
\definecolor{methcol}{RGB}{90,90,100}

\newcommand{\runtitle}{Thème 4 — Facile — Corrigé}
\pagestyle{fancy}\fancyhf{}
\fancyhead[L]{\small\textcolor{primarydark}{\textbf{Fiche 12} — Chimie organique}}
\fancyhead[R]{\small\textcolor{primarydark}{\runtitle}}
\fancyfoot[C]{\small\thepage}
\renewcommand{\headrulewidth}{0.8pt}
\renewcommand{\headrule}{\hbox to\headwidth{\color{primary}\leaders\hrule height \headrulewidth\hfill}}
\setlength{\parindent}{0pt}\setlength{\parskip}{3pt}

\tcbset{boxbase/.style={enhanced,breakable,boxrule=0.9pt,arc=2.5pt,
   left=3mm,right=3mm,top=2mm,bottom=2mm,fonttitle=\bfseries,coltitle=white,
   toptitle=0.5mm,bottomtitle=0.5mm}}
\newtcolorbox[auto counter]{corr}[1]{boxbase,colback=excol!5,colframe=excol,
   title={\textsc{Corrigé}~\thetcbcounter\ \textmd{--- \itshape #1}}}
\newcommand{\rep}[1]{\textcolor{primarydark}{\textbf{#1}}}

\newcommand{\banner}[2]{%
\begin{tcolorbox}[enhanced,colback=excol,colframe=excol,arc=5pt,boxrule=0pt,
   left=5mm,right=5mm,top=2.5mm,bottom=2.5mm,coltext=white]
{\large\bfseries #1}\\[1pt]{\small #2}
\end{tcolorbox}\vspace{2pt}}

\begin{document}

\banner{Thème 4 — Masse molaire d'une molécule organique}{Niveau \textbf{Facile} — Corrigé détaillé \quad$\vert$\quad Fiche 12, Chimie organique}

% ---------------------------------------------------------------------
\begin{corr}{le calcul modèle}
\[ M(\ce{C6H14N2}) = \underbrace{6\times12}_{72} + \underbrace{14\times1}_{14} + \underbrace{2\times14}_{28} = 72+14+28 = \rep{\qty{114}{\g\per\mol}}. \]
\end{corr}

% ---------------------------------------------------------------------
\begin{corr}{deux petites molécules}
\textbf{a)} $M(\ce{CH4}) = 12 + 4\times1 = \rep{\qty{16}{\g\per\mol}}$.\\
\textbf{b)} $M(\ce{C2H6O}) = 2\times12 + 6\times1 + 16 = 24+6+16 = \rep{\qty{46}{\g\per\mol}}$.
\end{corr}

% ---------------------------------------------------------------------
\begin{corr}{masse molaire ou masse moléculaire relative ?}
\textbf{a)} La masse molaire s'écrit avec son unité : $M = \rep{\qty{114}{\g\per\mol}}$.\\
\textbf{b)} La masse moléculaire relative est le même nombre \emph{sans unité} : $M_r = \rep{114}$.\\
\textbf{c)} $M_r$ est un \textbf{rapport} (à $\tfrac{1}{12}$ de la masse d'un atome de $\ce{^{12}C}$) : c'est une grandeur \rep{sans dimension}. Lui accoler «\,$\si{\g\per\mol}$\,» est donc \rep{faux} — cette unité est celle de $M$, pas de $M_r$.
\end{corr}

% ---------------------------------------------------------------------
\begin{corr}{de la structure à la masse molaire}
\textbf{a)} La propanone $\ce{CH3-CO-CH3}$ compte 3 C, 6 H et 1 O $\Rightarrow$ formule brute \rep{\ce{C3H6O}}.\\
\textbf{b)} $M = 3\times12 + 6\times1 + 16 = 36+6+16 = \rep{\qty{58}{\g\per\mol}}$.
\end{corr}

% ---------------------------------------------------------------------
\begin{corr}{de la masse à la quantité de matière}
\[ n = \frac{m}{M} = \frac{\qty{9,2}{\g}}{\qty{46}{\g\per\mol}} = \rep{\qty{0,20}{\mol}}. \]
(On aurait $N = n\,N_A = \qty{0,20}{}\times\qty{6,02e23}{} \approx \qty{1,2e23}{}$ molécules.)
\end{corr}

\end{document}
